\documentclass[11pt]{article}
\usepackage[T1]{fontenc}
\usepackage{lmodern}
\usepackage[margin=1in]{geometry}
\usepackage{amsmath,amssymb,amsthm,booktabs,hyperref}
\hypersetup{colorlinks=true,urlcolor=blue}
\newtheorem*{theorem}{Theorem}
\newcommand{\type}{\mathcal T}
\begin{document}
\begin{center}
{\Large\bfseries Identifying signed paths with a single probe}\\[5pt]
Akhil Ravuri \quad$\cdot$\quad Working note
\end{center}

Take two terminals joined by internally disjoint paths of lengths $2,1,3$.
Compare signs $(-,+,+)$ with $(-,-,-)$ in that order. Add a positive
path of length two to each graph:
\[
\begin{array}{c|rrrr}
\text{length}&2&1&3&2\\\midrule
\text{first graph}&-&+&+&+\\
\text{second graph}&-&-&-&+
\end{array}
\]
Negate the first row of signs, then exchange its two length-two paths.
The second row results. The observation defined below is invariant under
both operations. A positive probe of length two therefore fails on this
six-edge example. We will see why a probe of length four always works
on six edges.

\paragraph{The observation.}
Give each edge $uv$ a sign $b_{uv}\in\{1,-1\}$. Independently give each
vertex a uniform sign $U_v$. Count a conflict when
$U_u U_v=b_{uv}$. A path sign is the product of its edge signs.
The unknown type is the multiset of (length, path sign) pairs; the
individual edge signs along a path are not part of the target.
Add one positive path between the terminals, with fresh internal vertices.
Observe only the exact law of the total conflict count $X$ on this augmented
graph. Negating the vertex color at one terminal preserves its uniform law
and switches the effective signs of all paths.

\begin{theorem}
Suppose the original total length $m\ge6$ is known, there are at least
three paths, and at most one has length one. Among probe lengths $p\ge2$,
the shortest that identifies every original signed type is $m-2$.
\end{theorem}

\paragraph{What the law encodes.}
For an original path $i$ write $l_i$ for its length and $s_i$ for its sign.
Let $E(x)$ sum $\prod_{i\in I}s_i\,x^{\sum_{i\in I}l_i}$ over even-cardinality
subsets $I$ of paths (including the empty set); define $O(x)$ using odd
subsets. With $N=m+p$ and $G(z)=\mathbb E[z^X]$, the observable is
equivalently
\[
 C(x)=(1-x)^N G\!\left(\frac{1+x}{1-x}\right)=E(x)+x^p O(x).
\]
Indeed each edge contributes $1+x b_{uv}U_uU_v$ after normalization.
Averaging their product kills exactly the terms having an odd vertex degree.
Internal vertices thus require whole paths and the terminals require an
even number of paths. Including the probe leaves an odd original subset.
This derivation also shows that redistributing edge signs along a path
without changing their product preserves the law.

\paragraph{Return to the example.}
For the two original triples, subtract the second polynomial from the first:
$\Delta E=-2x^3-2x^5$ and $\Delta O=2x+2x^3$. At $p=2$ the terms cancel
exactly. At $p=4$ their sum $\Delta E+x^4\Delta O=-2x^3+2x^7$ is
nonzero: this probe tells the two graphs apart. The next argument shows
it works for \emph{every} six-edge pair, not just this one.

\newpage
\paragraph{Why $m-2$ suffices.}
Take two types giving the same $C$ and write $\Delta$ for their difference.
Since $\deg E\le m$ and $O(0)=0$, the only ambiguity when $p=m-2$ is
\[
 \Delta O=ax+bx^2,\qquad \Delta E=-a x^{m-1}-b x^m.
\]
There is at most one unit path, so $|a|\le2$. The coefficient of
$x^{m+p}$ in $C$ is $O_m$: it vanishes for an even number of paths and
is the product $\sigma$ of all path signs otherwise. Thus the two types
agree on the parity of their path counts.

If both counts are odd then $E_m=0$, forcing $b=0$. If both are even,
$E_{m-1}=\sigma E_1=0$: complementation preserves even cardinality, and
no even subset has length one. Hence $a=0$ and $|b|\le2$, since
$E_m=\sigma$. In either case at most one coefficient survives. But
\[
 (E-O)(1)=\prod_i(1-s_i)
\]
is divisible by eight for each type, since there are at least three paths.
The difference is $-2(a+b)$, of absolute value at most four. It must be
zero, so $E$ and $O$ agree separately.

To recover the paths, form
\[
 A=E+O=\prod_i(1+s_i x^{l_i}),\qquad
 D=(E+O)(E-O)=\prod_i(1-x^{2l_i}).
\]
At the shortest remaining length $l$ (half the first positive degree in
$D$), $-[x^{2l}]D$ counts its multiplicity, while $[x^l]A$ gives its
positive count minus its negative count. Divide both products by the
identified factors and repeat until $D=1$.

\paragraph{Why shorter probes fail.}
For each $2\le p\le m-3$, take lengths $(p,1,m-p-1)$ with signs
$(-,+,+)$ and $(-,-,-)$. Both graphs are simple and their original
types differ. After adding the positive length-$p$ path, the same negation
and exchange as in the opening example identifies the augmented laws.
This supplies a counterexample at every shorter admissible length.\hfill$\square$

\paragraph{Checks and context.}
A separate \href{https://github.com/Akhil-rav/discrete-math-lab/blob/0fb11c1/projects/finite-counting/check_single_probe.py}{enumeration}
conditions on the two terminal colors. At each fixed $m=6,\ldots,16$
all types give distinct laws for $p=m-2$ (4,164 types in total).
The normalized transform is binary coset Fourier duality: see A.~Goodall,
\href{https://iuuk.mff.cuni.cz/~andrew/dwajgApr06.pdf}{\emph{Fourier analysis on finite Abelian groups}},
eq.~(8), with the cut space as code and the edge signs as shift.
For the classical forward partition function on generalized theta graphs,
see A.~D.~Sokal, \href{https://arxiv.org/abs/cond-mat/0012369}{\emph{Chromatic
roots are dense in the whole complex plane}}, \S2.3. The question here is
inverse and assumes the host class and its original total length are known.
\end{document}
